\begin{register}{H}{AES\_KEY\_\regindex{n}\_REG (\regindex{n}: 0-7)}{\hexprefix{}0000+4*\regindex{n}}\label{regdesc:AESKEYnREG}
\regfield{}{32}{0}{{\hexprefix{}000000000}}%
\reglabel{Reset}\regnewline%
\begin{regdesc}\begin{reglist}
\item [AES\_KEY\_\regindex{n}\_REG (\regindex{n}: 0-7)]AES 密钥寄存器。（读/写）
\end{reglist}\end{regdesc}
\end{register}

\begin{register}{H}{AES\_TEXT\_IN\_\regindex{m}\_REG (\regindex{m}: 0-3)}{\hexprefix{}0020+4*\regindex{m}}\label{regdesc:AESTEXTINnREG}
\regfield{}{32}{0}{{\hexprefix{}000000000}}%
\reglabel{Reset}\regnewline%
\begin{regdesc}\begin{reglist}
\item [AES\_TEXT\_IN\_\regindex{m}\_REG (\regindex{m}: 0-3)] Typical AES 文本输入寄存器。（读/写）
\end{reglist}\end{regdesc}
\end{register}


\begin{register}{H}{AES\_TEXT\_OUT\_\regindex{m}\_REG (\regindex{m}: 0-3)}{\hexprefix{}0030+4*\regindex{m}}\label{regdesc:AESTEXTOUTnREG}
\regfield{}{32}{0}{{\hexprefix{}000000000}}%
\reglabel{Reset}\regnewline%
\begin{regdesc}\begin{reglist}
\item [AES\_TEXT\_OUT\_\regindex{m}\_REG (\regindex{m}: 0-3)] Typical AES 文本输出寄存器。（只读）
\end{reglist}\end{regdesc}
\end{register}


\begin{register}{H}{AES\_MODE\_REG}{\hexprefix{}0040}\label{regdesc:AESMODEREG}
\regfield{(reserved)}{29}{3}{{\hexprefix{}00000000}}%
\regfield{AES\_MODE}{3}{0}{{0}}%
\reglabel{Reset}\regnewline%
\begin{regdesc}\begin{reglist}
\item [AES\_MODE] 选择 AES 加速器的工作模式。详情请见表 \ref{table:AESOperation}。（读/写）
\end{reglist}\end{regdesc}
\end{register}


\begin{register}{H}{AES\_ENDIAN\_REG}{\hexprefix{}0044}\label{regdesc:AESENDIANREG}
\regfield{(reserved)}{26}{6}{{\hexprefix{}0000000}}%
\regfield{AES\_ENDIAN}{6}{0}{000000}%
\reglabel{Reset}\regnewline%
\begin{regdesc}\begin{reglist}
\item [AES\_ENDIAN] 字节序选择寄存器。详情请见表 \ref{table:TypicalAESPlaintextCiphertextEndian}。（读/写）
\end{reglist}\end{regdesc}
\end{register}


\begin{register}{H}{AES\_DMA\_ENABLE\_REG}{\hexprefix{}0090}\label{regdesc:AESDMAENABLEREG}
\regfield{(reserved)}{31}{1}{{\hexprefix{}00000000}}%
\regfield{AES\_DMA\_ENABLE}{1}{0}{{0}}%
\reglabel{Reset}\regnewline%
\begin{regdesc}\begin{reglist}
\item [AES\_DMA\_ENABLE] 选择工作模式。详情请见表 \ref{table:AES_DMA_ENABLE}。（读/写）
\end{reglist}\end{regdesc}
\end{register}


\begin{register}{H}{AES\_BLOCK\_MODE\_REG}{\hexprefix{}0094}\label{regdesc:AESBLOCKMODEREG}
\regfield{(reserved)}{29}{3}{{\hexprefix{}00000000}}%
\regfield{AES\_BLOCK\_MODE}{3}{0}{{0}}%
\reglabel{Reset}\regnewline%
\begin{regdesc}\begin{reglist}
\item [AES\_BLOCK\_MODE] 选择 DMA-AES 使用的块模式。详情请见表 \ref{table:AESBlockMode}。（读/写）
\end{reglist}\end{regdesc}
\end{register}


\begin{register}{H}{AES\_BLOCK\_NUM\_REG}{\hexprefix{}0098}\label{regdesc:AESBLOCKNUMREG}
\regfield{AES\_BLOCK\_NUM}{32}{0}{{\hexprefix{}00000000}}%
\reglabel{Reset}\regnewline%
\begin{regdesc}\begin{reglist}
\item [AES\_BLOCK\_NUM] 在 DMA-AES 运算中待加解密的文本块数。详情请见章节 \ref{Label Block Number}。（读/写）
\end{reglist}\end{regdesc}
\end{register}


\begin{register}{H}{AES\_INC\_SEL\_REG}{\hexprefix{}009C}\label{regdesc:AESINCSELREG}
\regfield{(reserved)}{31}{1}{{\hexprefix{}00000000}}%
\regfield{AES\_INC\_SEL}{1}{0}{{0}}%
\reglabel{Reset}\regnewline%
\begin{regdesc}\begin{reglist}
\item [AES\_INC\_SEL] 选择 CTR 块模式使用的标准增量函数。置 0 选择 INC$_{32}$ 标准增量函数，置 1 选择 INC$_{128}$ 标准增量函数。（读/写）
\end{reglist}\end{regdesc}
\end{register}


\begin{register}{H}{AES\_AAD\_BLOCK\_NUM\_REG}{\hexprefix{}00A0}\label{regdesc:AESAADBLOCKNUMREG}
\regfield{}{32}{0}{{\hexprefix{}00000000}}%
\reglabel{Reset}\regnewline%
\begin{regdesc}\begin{reglist}
\item [AES\_AAD\_BLOCK\_NUM] 在 GCM 运算中 AAD 的块数。详情请见章节 \ref{Label AAD Block Number}。（读/写）
\end{reglist}\end{regdesc}
\end{register}


\begin{register}{H}{AES\_REMAINDER\_BIT\_NUM\_REG}{\hexprefix{}00A4}\label{regdesc:AESREMAINDERBITNUMREG}
\regfield{(reserved)}{25}{7}{{\hexprefix{}0000000}}%
\regfield{AES\_REMAINDER\_BIT\_NUM}{7}{0}{{0}}%
\reglabel{Reset}\regnewline%
\begin{regdesc}\begin{reglist}
\item [AES\_REMAINDER\_BIT\_NUM] 在 GCM 运算中输入文本的不完整块的有效比特数。详情请见章节 \ref{Label Remainder Bit Number}。（读/写）
\end{reglist}\end{regdesc}
\end{register}



\begin{register}{H}{AES\_TRIGGER\_REG}{\hexprefix{}0048}\label{regdesc:AESTRIGGERREG}
\regfield{(reserved)}{31}{1}{{\hexprefix{}00000000}}%
\regfield{AES\_TRIGGER}{1}{0}{x}%
\reglabel{Reset}\regnewline%
\begin{regdesc}\begin{reglist}
\item [AES\_TRIGGER] 写入 1 使能 AES 运算。（只写）
\end{reglist}\end{regdesc}
\end{register}


\begin{register}{H}{AES\_STATE\_REG}{\hexprefix{}004C}\label{regdesc:AESSTATEREG}
\regfield{(reserved)}{30}{2}{{\hexprefix{}00000000}}%
\regfield{AES\_STATE}{2}{0}{{\hexprefix{}0}}%
\reglabel{Reset}\regnewline%
\begin{regdesc}\begin{reglist}
\item [AES\_STATE] AES 状态寄存器。详见表 \ref{table:AESStateValue-aes}（Typical AES 工作模式）和表 \ref{table:AESStateValue-dma}（DMA-AES 工作模式）。（只读）
\end{reglist}\end{regdesc}
\end{register}


\begin{register}{H}{AES\_CONTINUE\_OP\_REG}{\hexprefix{}00A8}\label{regdesc:AESCONTINUEOPREG}
\regfield{(reserved)}{31}{1}{{\hexprefix{}00000000}}%
\regfield{AES\_CONTINUE}{1}{0}{x}%
\reglabel{Reset}\regnewline%
\begin{regdesc}\begin{reglist}
\item [AES\_CONTINUE\_OP] 在 GCM 运算中，写入 1 继续运算。（只写）
\end{reglist}\end{regdesc}
\end{register}


\begin{register}{H}{AES\_DMA\_EXIT\_REG}{\hexprefix{}00B8}\label{regdesc:AESDMAEXITREG}
\regfield{(reserved)}{31}{1}{{\hexprefix{}00000000}}%
\regfield{AES\_DMA\_EXIT}{1}{0}{x}%
\reglabel{Reset}\regnewline%
\begin{regdesc}\begin{reglist}
\item [AES\_DMA\_EXIT] 在 DMA-AES 运算完成后，在下一次配置 AES 任何寄存器之前，写入 1 使 AES 回到空闲状态。（只写）
\end{reglist}\end{regdesc}
\end{register}


\begin{register}{H}{AES\_INT\_CLR\_REG}{\hexprefix{}00AC}\label{regdesc:AESINTCLRREG}
\regfield{(reserved)}{31}{1}{{\hexprefix{}00000000}}%
\regfield{AES\_INT\_CLR}{1}{0}{0}%
\reglabel{Reset}\regnewline%
\begin{regdesc}\begin{reglist}
\item [AES\_INT\_CLR] 写入 1 清除 AES 中断。（只写）
\end{reglist}\end{regdesc}
\end{register}


\begin{register}{H}{AES\_INT\_ENA\_REG}{\hexprefix{}00B0}\label{regdesc:AESINTENAREG}
\regfield{(reserved)}{31}{1}{{\hexprefix{}00000000}}%
\regfield{AES\_INT\_ENA}{1}{0}{0}%
\reglabel{Reset}\regnewline%
\begin{regdesc}\begin{reglist}
\item [AES\_INT\_ENA] 写入 1 使能 AES 中断功能，写入 0 关闭 AES 中断功能。（读/写）
\end{reglist}\end{regdesc}
\end{register}